\subsection{Weight-tail ideals: linear projective complexity and PD principality}
\label{arch136:tail-pdim}

\paragraph{Scope and status.}
These are explicit calculations for this investigation; no historical
priority is asserted.
The following two statements concern the coefficient algebra itself.
They constrain a bounded free \(S\)-linear resolution of all coefficient-tail
relations. They do not obstruct a nonlinear divided-power construction, and
they do not decide connectedness of \(\mathcal A_{\mathbb Z}\).
The first statement has been read in full by the author and independently
by the cote~15 reviewer (GPT-6 Astra Pro; effort not exposed;
Europe/Berlin date 2026-11-22). Root's whole-proof PASS is recorded
separately. The second statement is an elementary clarification of the
different behavior of ordinary modules and PD ideals.

\paragraph{ARCH-136-TAIL-PDIM.}
Let \(R=\mathbb Z\) or \(\mathbb Z_{(p)}\), let
\(S_R=\Gamma_R(T)\), with \(\operatorname{wt}(T)=1\), and let
\[
 I_{\ge N}=\bigoplus_{w\ge N}R\,T^{[w]},\qquad N\ge1.
\]
Then both \(I_{\ge N}\) and \(S_R/I_{\ge N}\) have infinite projective
dimension as ordinary \(S_R\)-modules. Projectives of arbitrary rank
are allowed in this assertion.

\emph{Proof.}
First work over \(k=\mathbb F_p\). Put \(x_i=T^{[p^i]}\). There is an
isomorphism of weight-graded \(k\)-algebras
\[
 \bar S=\Gamma_k(T)\ \cong
 \bigotimes_{i\ge0}k[x_i]/(x_i^p).
 \tag{TAIL.1}
\]
For completeness, if \(w=\sum_i e_i p^i\), \(0\le e_i<p\), the product
\(\prod_i x_i^{e_i}\) is a unit multiple of \(T^{[w]}\): its integer
coefficient is \(w!/\prod_i(p^i!)^{e_i}\), whose \(p\)-adic valuation is
zero by Legendre's formula. The coefficient of \(x_i^p\) has valuation
one, so \(x_i^p=0\) after reduction. These monomials give exactly one
basis vector in every weight. This is also the classical decomposition
in Cartan, \emph{Puissances divis\'ees}, Expos\'e~7, \S7, p.~7-10.

Choose \(a\ge0\) with \(p^a\ge N\), and set
\(D=k[x_a]/(x_a^p)\). Formula (TAIL.1) makes \(\bar S\) free as a
\(D\)-module. The quotient
\(\bar Q=\bar S/\bar I_{\ge N}\) has basis
\(T^{[0]},\ldots,T^{[N-1]}\), and \(x_a\) acts by zero, since it raises
weight by \(p^a\ge N\). Thus, as a \(D\)-module,
\(\bar Q\cong k^{\oplus N}\).

The standard resolution
\[
 \cdots\xrightarrow{x_a}D\xrightarrow{x_a^{p-1}}D
 \xrightarrow{x_a}D\longrightarrow k\longrightarrow0
 \tag{TAIL.2}
\]
is exact: the annihilators of \(x_a\) and \(x_a^{p-1}\) are respectively
\((x_a^{p-1})\) and \((x_a)\). Tensoring (TAIL.2) with \(k\) makes every
differential zero, so
\(\operatorname{Tor}^D_i(k,k)\cong k\) for every \(i\ge0\).
Consequently \(\bar Q\) has infinite \(D\)-projective dimension.
Restriction of scalars carries every \(\bar S\)-projective to a
\(D\)-projective, because \(\bar S\) is \(D\)-free. A finite
\(\bar S\)-projective resolution of \(\bar Q\) is therefore impossible.

Now let \(R=\mathbb Z_{(p)}\), or \(R=\mathbb Z\) and choose any prime
\(p\). If \(Q=S_R/I_{\ge N}\) had a finite \(S_R\)-projective resolution,
every term would be \(R\)-flat. The module \(Q\) is \(R\)-free.
Inductively the syzygies are \(R\)-flat, so reduction modulo \(p\)
preserves exactness of this resolution. Its terms remain projective
over \(S_R/pS_R\), contradicting the preceding paragraph.
Finally, the exact sequence
\(0\to I_{\ge N}\to S_R\to Q\to0\) shows that finite projective
dimension of \(I_{\ge N}\) would imply finite projective dimension of
\(Q\). This proves both assertions. \(\square\)

\paragraph{ARCH-136-PD-TAIL-PRINCIPAL: clarification at a prime.}
Let \(R=\mathbb Z_{(p)}\), \(N=p^a\), \(a\ge0\), and let \(u=T^{[N]}\).
Give the free PD generator \(U\) weight \(N\).
The PD homomorphism
\[
 \Gamma_R(U)\longrightarrow S_R,\qquad U\longmapsto u
 \tag{TAIL.3}
\]
is injective, and \(S_R\) is free over \(\Gamma_R(U)\), with basis
\[
 1,T^{[1]},\ldots,T^{[N-1]}.
\]
Moreover,
\[
 I_{\ge N}=\Gamma_R(U)_+S_R=(u)_{\mathrm{PD}},\qquad
 S_R/I_{\ge N}=S_R\otimes_{\Gamma_R(U)}R.
 \tag{TAIL.4}
\]
Thus this tail is PD-principal although its ordinary module projective
dimension is infinite.

\emph{Proof.}
Write \(w=rN+s\), \(0\le s<N\). The image of
\(\gamma_r(U)T^{[s]}\) is
\[
 \alpha_r\binom{rN+s}{s}\,T^{[rN+s]},\qquad
 \alpha_r=\frac{(rN)!}{r!(N!)^r}.
 \tag{TAIL.5}
\]
Legendre's formula and \(s_p(rN)=s_p(r)\) give
\(v_p(\alpha_r)=0\). Addition of \(s<N\) to \(rN\) has no base-\(p\)
carry, so the binomial coefficient is also a \(p\)-unit.
Formula (TAIL.5) therefore gives an \(R\)-basis, proving the freeness
and injectivity. The terms with \(r\ge1\) are exactly the weight tail,
and they lie in the PD ideal generated by \(u\); conversely that ideal
has no weights below \(N\). This proves (TAIL.4). \(\square\)

\paragraph{The exact limitation for the forward repair program.}
A repair that asks for a bounded projective \(S\)-resolution of
\(I_{\ge N}\) cannot use such a resolution: none exists.
This is a restriction on that linear implementation of simultaneous
coefficient-tail repairs. It does not say that every PD repair is
linear, nor that a source in \(\mathcal A_{\mathbb Z}\) must have
bounded module projective dimension. In particular (TAIL.4) shows why
PD principality and ordinary projective dimension must not be
conflated.

\paragraph{ARCH-136-POSITIVE-TAIL: coefficient PD generation is not positive-form generation.}
Let \(R=\mathbb Z_{(p)}\), \(S=\Gamma_R(T)\), and
\(N=p^a\), \(a\geq0\). Take a generator \(z\) of positive even
form degree \(q\geq2\) and positive weight \(m\). In the strict
PD square-zero extension
\[
             A=S\oplus Sz,
       \qquad z^2=0,\qquad \gamma_j(z)=0\quad(j\geq2),
 \tag{TAIL.6}
\]
with zero differential, the PD ideal generated by \(T^{[N]}z\)
has form-\(q\) part exactly
\[
                         S\,T^{[N]}z.
 \tag{TAIL.7}
\]
It does not contain \(T^{[pN]}z\), although the coefficient
tail \(I_{\ge N}\subset S\) is PD-principal by (TAIL.4).
The same non-generation holds after replacing \(Sz\) by
\((S/pS)z\).

\emph{Proof.}
The square-zero extension (TAIL.6) is the special PD extension:
on its square-zero ideal all higher divided powers are zero,
and mixed sums are defined by the divided-power addition identity.
Thus the ideal in question is already the ordinary ideal
\(S T^{[N]}z\). More generally, an operation of order \(j\geq2\)
on a homogeneous form-\(q\) generator has form degree \(jq>q\),
so it cannot provide additional form-\(q\) generators in a
nonnegative form grading.
At the required coefficient weight the only possible ordinary
product is a scalar multiple of
\[
 T^{[pN-N]}T^{[N]}z
       =\binom{pN}{N}T^{[pN]}z.
 \tag{TAIL.8}
\]
The binomial coefficient has \(p\)-adic valuation one. Hence it
cannot generate \(T^{[pN]}z\) over \(R\), and the product is
zero after replacing the module by \((S/pS)z\), whereas the
desired element is nonzero. In contrast,
\(\gamma_p(T^{[N]})\) is a \(p\)-unit times \(T^{[pN]}\),
by (TAIL.5). This proves the distinction. \(\square\)

This observation applies to the proposed families of primitives
of a positive-form class: all coefficients \(T^{[s]}z\) retain
the same form degree. A short resolution of the special conormal
\(I/I^{[2]}\) therefore cannot simply be substituted for a
resolution or a cycle lift of the module \(I z\). This is a
precise failure of that inference, not an obstruction to every
nonfree or nonlinear PD construction.

\paragraph{Source and operator scope.}
Cartan's cited decomposition is in
\url{https://www.numdam.org/item/SHC_1954-1955__7_1_A7_0.pdf}.
For a modern description of PD Beck modules, see A.~Kmail,
J.~Kozak and H.~Miller, \emph{Divided powers and K\"ahler differentials},
Proposition~4.2, Theorem~5.4 and Example~6.3, author version
\url{https://math.mit.edu/~hrm/papers/dp5.pdf}.
Its Beck differentials carry additional prime-power operators.
The special identity used here,
\(d\gamma_r(x)=\gamma_{r-1}(x)\,dx\), kills those additional
directions. Consequently a theorem about the Beck cotangent complex
does not by itself supply the strict special dg-PD repair or its
required uniform horizontal bounds. Dokas--Frankland--Ikonicoff,
\emph{Quillen (co)homology of divided power algebras over an operad},
arXiv:2403.18049v3, Theorem~8.2, constructs a cotangent complex in a
setting whose ground ring is assumed to be a field (p.~4).

